\documentclass[11pt,a4paper]{article}
\usepackage[margin=1in]{geometry}
\usepackage{amsmath,amssymb,amsthm}
\usepackage{algorithm}
\usepackage{algpseudocode}
\usepackage{booktabs}
\usepackage{hyperref}

\theoremstyle{definition}
\newtheorem{definition}{Definition}[section]
\newtheorem{theorem}{Theorem}[section]
\newtheorem{lemma}[theorem]{Lemma}
\newtheorem{remark}{Remark}[section]

\title{\textbf{Rigorous Mathematical Specification and Memory Caching of Segment-Level Recurrence}}
\author{Team Scaibu \\ \small Email: \texttt{jaydeep.9664920749@gmail.com} \\ \small Architecture and Core Engine Team}
\date{October 2026}

\begin{document}
\maketitle

\begin{abstract}
Fixed-length context limits induce context fragmentation. Segment-level recurrence (Transformer-XL) caches previous segment hidden states as non-differentiable historical memory $\mathbf{M}$, concatenating them along the sequence axis to allow attention across multiple segment horizons without backpropagation through memory.
\end{abstract}

\section{Mathematical Formulation}

Let current segment be $\mathbf{S}_\tau \in \mathbb{R}^{L \times d}$ and cached memory be $\mathbf{M}_{\tau-1} \in \mathbb{R}^{M \times d}$. Extended representation is:
\begin{equation}
\tilde{\mathbf{H}}_\tau = [\operatorname{stop\_gradient}(\mathbf{M}_{\tau-1}); \mathbf{S}_\tau] \in \mathbb{R}^{(M + L) \times d}
\end{equation}
Queries attend across $\tilde{\mathbf{H}}_\tau$, enabling information reuse across segments.

\begin{thebibliography}{9}
\bibitem{dai2019} Z.~Dai et al., ``Transformer-XL: Attentive language models beyond a fixed-length context,'' in \emph{ACL}, 2019.
\end{thebibliography}

\end{document}
